\documentclass[11pt]{article}
\usepackage[margin=1in]{geometry}
\usepackage{amsmath,amssymb}
\usepackage{booktabs}
\usepackage{array}
\usepackage{xcolor}
\usepackage{caption}
\usepackage{enumitem}
\usepackage[hidelinks]{hyperref}
\captionsetup{font=small,labelfont=bf}
\setlength{\parskip}{4pt}
\setlength{\parindent}{0pt}
\emergencystretch=1.5em

\newcommand{\champ}[1]{\textbf{#1}\,$^{\bigstar}$}

\title{\textbf{R-Value: Player and Team Production per Dollar\\ of True NBA Expenditure}\\[6pt]
\large Model definition, six-season backtest (2018--19 to 2025--26),\\
and 2026--27 projections}
\author{NBA Analytics Project}
\date{July 5, 2026}

\begin{document}
\maketitle

\begin{abstract}
\noindent
R-Value measures NBA production per dollar of \emph{true} team expenditure ---
salary grossed up by each team's luxury-tax burden under the CBA schedule in
force that season --- normalized so the league average is 100. Its performance
core is a calibrated blend of five statistical pillars; its cost core prices
every contract against the era-correct tax structure (2017-CBA brackets through
2022--23, the 2023-CBA regime thereafter). Calibrated on six seasons ---
2018--19 and 2021--22 through 2025--26, deliberately excluding the two
COVID-distorted years --- the performance core identifies All-NBA selections
better than PIE, raw plus-minus, or estimated net rating, and a fitted title
score ranked the actual champion \#1, \#2, \#6, \#1, \#2, \#3. Under
leave-one-season-out cross-validation, in which every fitted parameter is
refit without the season it is scored on, it holds \textbf{recall@15 $=
0.789$} (vs.\ 0.711 for the best unfitted rival), \textbf{$\rho = 0.914$}
against team wins, and a mean champion rank of \textbf{2.83} of 30 teams. For
2026--27 the model projects Oklahoma City as champion, New York and Houston as
\#1 seeds, and Detroit as the league's most cost-efficient team.
\end{abstract}

\section{Key findings}
\begin{itemize}[itemsep=2pt]
\item \textbf{The performance core is the best award predictor tested, and it
survives cross-validation.} Over six seasons it recovers $82.2\%$ of actual
All-NBA selections in its top 15, versus $71.1\%$ for PIE$\times$minutes,
$34.4\%$ for raw plus-minus, and $23.3\%$ for estimated net rating. Refit
per fold without the season it is scored on, it still recovers $78.9\%$ ---
the gap between fitted and honest accuracy is three percentage points
(Section~\ref{sec:rivals}).
\item \textbf{Extending the calibration window rebalanced the weights.} On
three seasons the fit leaned on box production (0.35) and the impact ratings
(0.20 offense, 0.20 defense). On six seasons the offensive and defensive
impact weights each settle at \textbf{0.15}, box drops to 0.25, and the freed
weight flows to PIE (0.30) and load-scaled efficiency (0.15): over a longer
window, noisy on-court impact ratings cede ground to the steadier all-in-one
and efficiency signals (Section~\ref{sec:calib}).
\item \textbf{Rotation depth, not star power, separates champions.} Fitted on
six actual champions, the title score weights the eight-man-rotation term nine
times heavier than the star term ($a = 0.2$ on the top-two, $b = 1.8$ on the
rotation). Star ceiling matters, but the quality of the de facto playoff
rotation is what the fit keeps choosing (Section~\ref{sec:backtest}).
\item \textbf{Champions usually are not value teams.} Across the six title
winners, value ranks were 16th, 29th, 14th, 13th, \textbf{3rd}, and 14th. The
2021--22 Warriors are the extreme: a \$179M payroll plus a \$170M repeater-tax
bill bought a championship at a Team R-Value of 49 --- the least efficient
title in the sample. Only the 2024--25 Thunder won while near the top of the
value board, the signature of a rookie-scale core peaking before its
extensions kick in.
\item \textbf{2026--27:} Oklahoma City is a clear title favorite (score 4.49
vs.\ New York 3.80, Houston 3.63, Denver 3.42) despite projecting as only the
West \#4 seed on per-minute quality --- the league's best top-two plus its
deepest rotation, now paid for in full (\$234M payroll, \$108M tax). Detroit
projects as the most R-Valuable team (159) for the second straight season.
\end{itemize}

\section{The model}
\label{sec:model}

\subsection{Performance pillars}
For each player-season, five pillars are computed and standardized. Let
$z(\cdot)$ denote a z-score whose mean and deviation are taken over the
qualified pool (players with $m \ge 500$ minutes), and let ``lg'' denote the
minutes-weighted league average. With $m$ the player's minutes:
\begin{align}
\mathrm{BOX} &= \mathrm{PTS} + 0.7\,\mathrm{REB} + 1.4\,\mathrm{AST}
  + 2.2\,\mathrm{STL} + 2.0\,\mathrm{BLK} - 1.4\,\mathrm{TOV} \\[2pt]
z_{\mathrm{box}} &= \tfrac12\, z\!\left(36\,\mathrm{BOX}/m\right)
  + \tfrac12\, z\!\left(\mathrm{BOX}\right)
  && \text{(rate + volume)} \\[2pt]
z_{\mathrm{eff}} &= z\!\left[(\mathrm{TS\%} - \mathrm{TS\%}^{\mathrm{lg}})
  \cdot \mathrm{USG\%}\right]
  && \text{(load-scaled efficiency)} \\[2pt]
z_{\mathrm{off}} &= z\!\left[\bigl(0.6\,\Delta E_{\mathrm{OFF}}
  + 0.4\,\Delta \mathrm{OFF}\bigr)\cdot \tfrac{m}{m+500}\right] \\[2pt]
z_{\mathrm{def}} &= z\!\left[\bigl(0.6\,\Delta E_{\mathrm{DEF}}
  + 0.4\,\Delta \mathrm{DEF}\bigr)\cdot \tfrac{m}{m+500}\right] \\[2pt]
z_{\mathrm{adv}} &= z\!\left(\mathrm{PIE}\right)
  && \text{(game-event share)}
\end{align}
where $\Delta$ terms are the player's (estimated) offensive/defensive rating
relative to league average, and the $m/(m+500)$ factor is a reliability
shrinkage that pulls low-minute impact ratings toward zero.

\subsection{Composite score, availability, and WPS}
\begin{align}
\mathrm{CPI} &= 0.25\,z_{\mathrm{box}} + 0.15\,z_{\mathrm{eff}}
  + 0.15\,z_{\mathrm{off}} + 0.15\,z_{\mathrm{def}} + 0.30\,z_{\mathrm{adv}}
  && \text{(calibrated, Sec.~\ref{sec:calib})} \\[2pt]
\mathrm{PS} &= 100\,\Phi(\mathrm{CPI})
  && \text{($\Phi$ = standard normal CDF)} \\[2pt]
A &= \sqrt{\min(m,\,2200)/2200}
  && \text{(availability credit)} \\[2pt]
\mathrm{WPS} &= \mathrm{PS}\cdot\bigl(0.55 + 0.45\,A\bigr)
  && \text{(weighted performance score)}
\end{align}
PS maps CPI onto a 0--100 percentile-like scale; WPS discounts it for games
missed, with a floor so that an injured star is dampened, not erased.

\subsection{Cost side: pricing each era's CBA}
Team payroll $P$ is the sum of guaranteed salaries (dead money included,
two-way contracts excluded). The luxury tax $T(P)$ is computed from the exact
bracket schedule of the applicable league year, using repeater rates where a
team qualified:
\begin{equation}
T(P) \;=\; \sum_{k\ge 1} r_k \cdot
  \min\!\Bigl(w,\; \max\bigl(0,\; P - \tau - (k-1)\,w\bigr)\Bigr),
\end{equation}
with $\tau$ the tax line and $w$ the bracket width. Under the 2017 CBA
(through 2022--23): $w = \$5$M, rates 1.50, 1.75, 2.50, 3.25, \ldots{}
($+1.00$ each for repeaters). Under the 2023 CBA (2025--26 rates, width
indexed to cap growth): 1.00, 1.25, 3.50, 4.75, \ldots{} standard and 3.00,
3.25, 5.50, 6.75, \ldots{} repeater, $+0.50$ per bracket beyond the fourth.
Total expenditure and each player's effective cost are then
\begin{equation}
\mathrm{TE} = P + T(P), \qquad
\mathrm{EC}_i = s_i \cdot \frac{\mathrm{TE}}{P}, \qquad
\mathrm{CI}_i = \operatorname{clip}\!\left(
  \frac{\mathrm{EC}_i}{\overline{\mathrm{TE}}},\; 0.002,\; 0.40\right),
\end{equation}
where $s_i$ is the player's salary and $\overline{\mathrm{TE}}$ the league-mean
total expenditure that season. Every salary dollar on a taxpaying roster is
grossed up by that team's tax multiplier: a \$40M contract on a
repeater-taxpayer costs its franchise far more than \$40M.

\subsection{R-Value}
\begin{align}
\textbf{Player:}\quad
R_i &= 100 \cdot \operatorname{normalize}\!\left[
  \mathrm{WPS}_i \cdot \left(\frac{\overline{\mathrm{CI}}}{\mathrm{CI}_i}\right)^{0.35}
  \right] \\[4pt]
\textbf{Team:}\quad
\mathrm{TeamR} &= 100 \cdot \operatorname{normalize}\!\left[
  \frac{\sum_i \mathrm{WPS}_i\, m_i}{\mathrm{TE}} \right] \\[4pt]
\textbf{Title:}\quad
\mathrm{TITLE} &= z(\mathrm{PERF}) \;+\; a\, z(\mathrm{TOP2})
  \;+\; b\, z(\mathrm{ROT8}), \qquad a = 0.2,\;\; b = 1.8
\end{align}
where $\operatorname{normalize}$ divides by the league mean (so 100 is
average) and the $0.35$ exponent dampens cost leverage so minimum contracts
cannot trivially top the list.
$\mathrm{PERF} = \sum \mathrm{WPS}\,m / \sum m$ is minutes-weighted team
quality; $\mathrm{TOP2}$ is the mean WPS of the two best players; and
$\mathrm{ROT8}$ is the minutes-weighted WPS of the eight most-played
players --- the de facto rotation that playoff basketball compresses to.

\subsection{Calibration and data}
\label{sec:calib}
The five CPI weights are fitted by a two-stage simplex grid search maximizing
an equal blend of All-NBA recall@15 and Spearman correlation with team wins,
averaged over the six calibration seasons. The title premiums $(a,b)$ are then
fitted on a $0$--$2.5$ grid to minimize the six actual champions' mean
title-score rank, preferring smaller premiums on ties.

Two seasons inside the window are deliberately excluded: 2019--20 (the bubble
season; teams played uneven schedules, breaking the availability model) and
2020--21 (72 games; luxury-tax bills were pro-rated at settlement, breaking
the cost model). Award eligibility follows each era's rules: the 65-game
requirement applies from 2023--24 only, with a 40-game floor earlier.
Pre-2023 salaries are reconstructed from per-season contract archives,
individual player salary histories where re-signings purged the archive rows,
and hand-verified patches for players absent from all sources; team payroll
totals are validated against independent per-team aggregates for all 30 teams
in every season.

\section{Validation against rival metrics}
\label{sec:rivals}

\begin{table}[h]\centering
\caption{Predicting All-NBA selections: recall@15 among award-eligible
players, by season. R-Value uses its cost-free performance core (WPS).}
\small
\begin{tabular}{lccccccc}
\toprule
Metric & 18--19 & 21--22 & 22--23 & 23--24 & 24--25 & 25--26 & Mean \\
\midrule
\textbf{R-Value (WPS core)} & \textbf{0.800} & \textbf{0.867} & 0.733 & \textbf{0.933} & \textbf{0.800} & \textbf{0.800} & \textbf{0.822} \\
PIE $\times$ minutes        & 0.733 & 0.733 & 0.733 & 0.733 & 0.600 & 0.733 & 0.711 \\
Raw plus-minus              & 0.333 & 0.400 & 0.267 & 0.333 & 0.400 & 0.333 & 0.344 \\
Estimated net rating        & 0.200 & 0.267 & 0.200 & 0.200 & 0.267 & 0.267 & 0.233 \\
\bottomrule
\end{tabular}
\end{table}

\begin{table}[h]\centering
\caption{Predicting regular-season wins: Spearman rank correlation, by season.}
\small
\begin{tabular}{lccccccc}
\toprule
Metric & 18--19 & 21--22 & 22--23 & 23--24 & 24--25 & 25--26 & Mean \\
\midrule
Team net rating            & 0.976 & 0.950 & 0.931 & 0.955 & 0.952 & 0.951 & 0.952 \\
\textbf{R-Value perf core} & \textbf{0.935} & \textbf{0.902} & \textbf{0.880} & \textbf{0.924} & \textbf{0.932} & \textbf{0.927} & \textbf{0.917} \\
Team R-Value (per dollar)  & 0.531 & 0.482 & 0.297 & 0.499 & 0.633 & 0.813 & 0.543 \\
Payroll                    & 0.431 & 0.360 & 0.575 & 0.611 & 0.129 & 0.393 & 0.417 \\
\bottomrule
\end{tabular}
\end{table}

Three observations. First, R-Value beats or ties every rival on All-NBA
prediction in five of six seasons (2022--23 is a tie with PIE; voters that
year rewarded narrative picks like Fox and Randle that no impact metric
selects). Second, team net rating edges the performance core on wins ($0.952$
vs.\ $0.917$) --- but net rating \emph{is} point differential, which is
near-definitionally wins; it cannot rank players, price contracts, or be
projected forward from individual contracts. R-Value does all three from one
framework. Third, payroll alone correlates with wins at only $0.417$ ---
money buys far less than reputation suggests, which is precisely the
inefficiency R-Value quantifies.

\subsection{Out-of-sample validation}
\label{sec:cv}

The figures above are \emph{in-sample}: the pillar weights and title premiums
were grid-searched to maximize exactly those quantities on exactly those
seasons. The rival metrics have no fitted parameters, so this comparison
flatters R-Value. To correct for it, every fitted parameter is refit on five
seasons and scored on the sixth, rotating through all six.

\begin{table}[h]\centering
\caption{Leave-one-season-out cross-validation. Each row refits the five
pillar weights and both title premiums \emph{without} the held-out season,
then scores that season.}
\small
\begin{tabular}{lcccccccccc}
\toprule
Held out & Recall & $\rho$ & Champ & \multicolumn{5}{c}{Weights chosen} & $a$ & $b$ \\
\cmidrule(lr){5-9}
 & @15 & wins & rank & box & eff & off & def & PIE & & \\
\midrule
2018--19 & 0.800 & 0.935 & 1 & 0.25 & 0.15 & 0.15 & 0.15 & 0.30 & 0.2 & 1.8 \\
2021--22 & 0.800 & 0.882 & 3 & 0.40 & 0.05 & 0.15 & 0.20 & 0.20 & 1.5 & 2.4 \\
2022--23 & 0.733 & 0.880 & 6 & 0.25 & 0.15 & 0.15 & 0.15 & 0.30 & 0.2 & 1.8 \\
2023--24 & 0.800 & 0.926 & 1 & 0.25 & 0.20 & 0.15 & 0.15 & 0.25 & 0.1 & 2.1 \\
2024--25 & 0.800 & 0.932 & 3 & 0.25 & 0.15 & 0.15 & 0.15 & 0.30 & 0.0 & 1.8 \\
2025--26 & 0.800 & 0.927 & 3 & 0.25 & 0.15 & 0.15 & 0.15 & 0.30 & 0.2 & 1.8 \\
\midrule
\textbf{Mean} & \textbf{0.789} & \textbf{0.914} & \textbf{2.83} & \multicolumn{7}{c}{} \\
\bottomrule
\end{tabular}
\end{table}

The degradation is small: recall@15 falls from $0.822$ to $0.789$, still
comfortably ahead of PIE$\times$minutes at $0.711$; the wins correlation falls
from $0.917$ to $0.914$; mean champion rank slips from $2.50$ to $2.83$. The
weights are also stable --- four of six folds select an identical vector, the
offensive-impact weight is $0.15$ in \emph{all six}, and the defensive weight
varies by $0.05$. Only the 2021--22 fold diverges materially, leaning on box
production instead of PIE.

The title premiums tell the same story: the rotation term $b$ lands in
$[1.8, 2.4]$ in every fold while the star term $a$ sits at or below $0.2$ in
five of six. Whatever the precise magnitudes, the fit keeps choosing rotation
depth over star ceiling, and it does so without having seen the champion it is
asked to rank.

\section{Historical backtest: six seasons, six champions}
\label{sec:backtest}

\subsection{Did the title score find the champion?}

\begin{table}[h]\centering
\caption{Top five title scores per season vs.\ what actually happened.
$^{\bigstar}$ = actual champion. PO-W = playoff wins (16 = title).}
\scriptsize
\begin{tabular}{clcccccc}
\toprule
Season & Team & Title & Perf & TOP2 & ROT8 & Wins & PO-W \\
\midrule
2018--19 & \champ{TOR} & 5.12 & 68.2 & 92.5 & 76.8 & 58 & 16 \\
         & GSW         & 4.52 & 64.8 & 98.2 & 74.3 & 57 & 14 \\
         & MIL         & 4.38 & 62.9 & 94.1 & 75.1 & 60 & 10 \\
         & UTA         & 3.58 & 64.0 & 89.7 & 69.7 & 50 & 1  \\
         & PHI         & 3.39 & 61.3 & 92.4 & 69.9 & 51 & 7  \\
\midrule
2021--22 & BOS         & 5.09 & 67.1 & 92.9 & 74.7 & 51 & 14 \\
         & \champ{GSW} & 3.72 & 61.7 & 87.5 & 70.0 & 53 & 16 \\
         & PHX         & 3.72 & 62.6 & 92.2 & 69.0 & 64 & 7  \\
         & UTA         & 3.16 & 59.2 & 92.3 & 67.6 & 49 & 2  \\
         & MEM         & 2.99 & 60.3 & 85.2 & 66.5 & 56 & 6  \\
\midrule
2022--23 & BOS         & 4.45 & 65.4 & 90.4 & 71.0 & 57 & 11 \\
         & CLE         & 4.02 & 63.6 & 89.6 & 69.7 & 51 & 1  \\
         & SAC         & 2.81 & 59.7 & 90.4 & 65.1 & 48 & 3  \\
         & MIL         & 2.67 & 56.5 & 92.6 & 66.1 & 58 & 1  \\
         & NYK         & 2.40 & 56.7 & 84.5 & 65.2 & 47 & 6  \\
         & \multicolumn{7}{l}{\emph{champion DEN ranked \#6 (coasted to perf \#9; playoff Joki\'c broke the model)}} \\
\midrule
2023--24 & \champ{BOS} & 5.29 & 70.7 & 91.6 & 78.8 & 64 & 16 \\
         & MIN         & 3.98 & 68.5 & 90.1 & 71.1 & 56 & 9  \\
         & OKC         & 3.28 & 61.5 & 94.9 & 70.2 & 57 & 6  \\
         & NOP         & 2.92 & 61.3 & 83.8 & 68.9 & 49 & 0  \\
         & LAC         & 2.80 & 60.0 & 91.0 & 68.1 & 51 & 2  \\
\midrule
2024--25 & CLE         & 4.89 & 65.3 & 94.5 & 78.9 & 64 & 5  \\
         & \champ{OKC} & 4.32 & 67.1 & 94.5 & 73.5 & 68 & 16 \\
         & BOS         & 4.31 & 68.0 & 88.0 & 73.6 & 61 & 6  \\
         & MIN         & 3.20 & 63.3 & 86.1 & 68.4 & 49 & 9  \\
         & NYK         & 2.84 & 60.9 & 89.7 & 66.8 & 51 & 10 \\
\midrule
2025--26 & SAS         & 4.06 & 66.0 & 89.0 & 72.6 & 62 & 13 \\
         & BOS         & 3.55 & 65.0 & 88.0 & 69.6 & 56 & 3  \\
         & \champ{NYK} & 3.37 & 61.5 & 89.6 & 70.3 & 53 & 16 \\
         & OKC         & 3.18 & 62.0 & 97.1 & 67.7 & 64 & 11 \\
         & MIN         & 3.12 & 60.4 & 85.4 & 69.6 & 49 & 6  \\
\bottomrule
\end{tabular}
\end{table}

The champion ranked \textbf{\#1, \#2, \#6, \#1, \#2, \#3} --- mean rank 2.50
out of 30 teams, with four of six champions in the top two. The one real miss
is instructive: the 2022--23 Nuggets ranked ninth in regular-season
performance because Joki\'c's supporting cast coasted, then flipped a switch
in April --- exactly the failure mode a regular-season metric cannot see. The
2021--22 season shows the opposite face: Boston graded far ahead of the field,
reached the Finals, and lost it to the Warriors in six --- a near-hit ranked
second.

\subsection{Value and winning are different axes}

\begin{table}[h]\centering
\caption{Most R-Valuable teams per season vs.\ where the champion ranked in
value terms (of 30).}
\small
\begin{tabular}{clcc}
\toprule
Season & Top three by Team R-Value & Champion (TeamR) & Champ.\ value rank \\
\midrule
2018--19 & PHI 139, UTA 138, IND 136 & TOR (105) & 16th \\
2021--22 & MEM 152, CHA 149, PHX 144 & GSW (49)  & 29th \\
2022--23 & MEM 151, SAC 142, NYK 141 & DEN (107) & 14th \\
2023--24 & OKC 143, ORL 142, IND 139 & BOS (109) & 13th \\
2024--25 & DET 153, CLE 151, OKC 151 & OKC (151) & 3rd  \\
2025--26 & MIA 141, SAS 138, DET 136 & NYK (103) & 14th \\
\bottomrule
\end{tabular}
\end{table}

The pattern across eight years is stable: championships are bought at roughly
league-average efficiency or worse, and the value leaderboard is populated by
young teams on rookie-scale contracts --- some of which (OKC 2023--24
$\rightarrow$ champion 2024--25) convert value into titles, and some of which
(MEM 2021--23) never do. The 2021--22 Warriors bookend the distribution: with
a \$179M payroll and a \$170M repeater-tax bill, their Team R-Value of 49 is
the least efficient championship in the sample --- and a working definition of
what ``buying a title'' looks like in R-Value terms.

\subsection{All-NBA through the R-Value lens}
Of the 90 actual All-NBA selections across the six seasons, the performance
core placed \textbf{74 (82.2\%)} in its top fifteen eligible players
(12, 13, 11, 14, 12, and 12 per season). The same tables expose the price of
stardom in every era: in 2021--22, rookie-scale Luka Don\v{c}i\'c (WPS 96.9,
\$10.2M, R $= 210$) and Trae Young (WPS 95.1, \$8.3M, R $= 221$) delivered
first-five production at a fifth of the veteran max; in 2025--26 the same slot
belongs to Jalen Duren (WPS 95.7, \$6.5M, R $= 253$) and Victor Wembanyama
(WPS 96.2, \$13.4M, R $= 197$).

\section{2026--27 projections}
\label{sec:predict}

Player CPI is projected as a $0.55/0.30/0.15$ blend of the last three seasons
plus an age-curve adjustment; minutes from a two-season blend, rescaled to a
full 19{,}680-minute roster; rosters and salaries reflect contracts signed as
of July 5, 2026 (431 players, of whom 41 rookies/unknowns enter at replacement
level). Unsigned free agents are absent by construction.

\subsection{Predicted standings}

\begin{table}[h]\centering
\caption{Projected 2026--27 standings by performance core. TeamR = projected
Team R-Value; payroll and tax in \$M.}
\scriptsize
\begin{tabular}{clcccrr@{\hspace{16pt}}clcccrr}
\toprule
\multicolumn{7}{c}{\textbf{East}} & \multicolumn{7}{c}{\textbf{West}} \\
\cmidrule(r{16pt}){1-7}\cmidrule{8-14}
\# & Team & Perf & Title & TeamR & Pay & Tax &
\# & Team & Perf & Title & TeamR & Pay & Tax \\
\midrule
1  & NYK & 60.9 & 3.80  & 95  & 215 & 46 & 1  & HOU & 60.5 & 3.63  & 125 & 199 & 0   \\
2  & CLE & 55.4 & 2.40  & 123 & 184 & 0  & 2  & DEN & 59.2 & 3.42  & 83  & 220 & 73  \\
3  & BOS & 55.3 & 2.47  & 109 & 203 & 5  & 3  & SAS & 58.9 & 3.00  & 122 & 198 & 0   \\
4  & TOR & 55.0 & 2.29  & 111 & 202 & 1  & 4  & OKC & 58.0 & 4.49  & 69  & 234 & 108 \\
5  & DET & 54.4 & 2.67  & 159 & 140 & 0  & 5  & MIN & 56.5 & 2.22  & 80  & 219 & 70  \\
6  & ORL & 54.4 & 1.69  & 80  & 224 & 55 & 6  & LAL & 52.8 & 1.46  & 112 & 192 & 0   \\
7  & PHI & 51.6 & 0.86  & 103 & 203 & 2  & 7  & GSW & 52.7 & 0.90  & 117 & 184 & 0   \\
8  & IND & 51.3 & 0.54  & 98  & 208 & 7  & 8  & POR & 46.9 & $-$1.17 & 101 & 191 & 0 \\
9  & MIA & 49.5 & 0.18  & 102 & 199 & 0  & 9  & PHX & 45.6 & $-$0.87 & 73  & 214 & 43 \\
10 & ATL & 48.4 & 0.60  & 104 & 191 & 0  & 10 & DAL & 43.9 & $-$1.48 & 101 & 178 & 0  \\
11 & CHA & 47.2 & $-$0.01 & 115 & 169 & 0 & 11 & NOP & 42.6 & $-$1.99 & 91  & 192 & 0 \\
12 & MIL & 44.0 & $-$0.40 & 108 & 167 & 0 & 12 & LAC & 40.6 & $-$3.00 & 122 & 136 & 0 \\
13 & CHI & 41.7 & $-$2.14 & 104 & 165 & 0 & 13 & UTA & 37.6 & $-$4.37 & 88  & 175 & 0 \\
14 & BKN & 36.2 & $-$4.38 & 90  & 164 & 0 & 14 & SAC & 37.1 & $-$4.56 & 60  & 219 & 35 \\
15 & WAS & 32.5 & $-$6.76 & 73  & 183 & 0 & 15 & MEM & 33.1 & $-$5.49 & 86  & 158 & 0  \\
\bottomrule
\end{tabular}
\end{table}

\subsection{Title picture}

\begin{table}[h]\centering
\caption{Top five 2026--27 title scores. OKC wins on rotation depth despite
projecting fourth in the West on per-minute quality.}
\begin{tabular}{clccccc}
\toprule
Rank & Team & Conf.\ seed & Title score & Perf & TOP2 & ROT8 \\
\midrule
1 & \textbf{OKC} & West \#4 & \textbf{4.49} & 58.0 & 93.6 & \textbf{71.6} \\
2 & NYK & East \#1 & 3.80 & 60.9 & 88.6 & 66.2 \\
3 & HOU & West \#1 & 3.63 & 60.5 & 87.1 & 65.7 \\
4 & DEN & West \#2 & 3.42 & 59.2 & 92.9 & 64.8 \\
5 & SAS & West \#3 & 3.00 & 58.9 & 85.9 & 63.5 \\
\bottomrule
\end{tabular}
\end{table}

OKC's projected eight-man rotation is more than five WPS points clear of every
other contender --- the same profile the backtest ranked top-two in 2024--25
--- while Denver enters the contender group on the strength of the
Joki\'c--Murray top-two (92.9). The champion pick is robust to the calibration
window: both the three-season and six-season fits select OKC, but the
six-season fit widens the margin. The cost side is equally striking: OKC's
Team R-Value falls to 69 (\$234M payroll, \$108M tax), completing the
conversion from the league's best value team (2023--24) to its biggest
spender --- the exact path the 2021--22 Warriors took to their inefficient
title.

\subsection{Predicted All-NBA}

\begin{table}[h]\centering
\caption{Projected 2026--27 All-NBA: top 15 by projected WPS among players
projected for $\ge 1600$ minutes. Salary in \$M.}
\small
\begin{tabular}{clcccrr}
\toprule
Team & Player & Club & Age & WPS & Salary & R-Value \\
\midrule
1st & Shai Gilgeous-Alexander & OKC & 28 & 99.8 & 40.8 & 126 \\
1st & Nikola Joki\'c          & DEN & 32 & 99.8 & 59.0 & 114 \\
1st & Luka Don\v{c}i\'c       & LAL & 28 & 96.9 & 51.0 & 129 \\
1st & Victor Wembanyama       & SAS & 23 & 94.0 & 16.9 & 184 \\
1st & Jalen Duren             & DET & 23 & 92.9 &  9.0 & 227 \\
\midrule
2nd & Donovan Mitchell        & CLE & 30 & 91.9 & 50.1 & 123 \\
2nd & Karl-Anthony Towns      & NYK & 31 & 91.5 & 57.1 & 109 \\
2nd & Anthony Edwards         & MIN & 25 & 90.7 & 48.9 & 111 \\
2nd & Cade Cunningham         & DET & 25 & 90.5 & 50.1 & 121 \\
2nd & Tyrese Haliburton       & IND & 27 & 90.5 & 48.9 & 121 \\
\midrule
3rd & Kawhi Leonard           & TOR & 36 & 88.9 & 50.3 & 119 \\
3rd & Kevin Durant            & HOU & 38 & 87.8 & 43.9 & 123 \\
3rd & Chet Holmgren           & OKC & 25 & 87.3 & 41.2 & 110 \\
3rd & Evan Mobley             & CLE & 26 & 87.1 & 50.1 & 117 \\
3rd & Jarrett Allen           & CLE & 29 & 86.7 & 28.0 & 142 \\
\bottomrule
\end{tabular}
\end{table}

\subsection{Where the value is}

\begin{table}[h]\centering
\caption{Projected most R-Valuable teams (left) and highest player R-Values
among players projected for $\ge 1200$ minutes (right, top 8).}
\small
\begin{tabular}{clcc@{\hspace{24pt}}clcc}
\toprule
\# & Team & TeamR & Proj.\ seed & \# & Player & R & WPS \\
\midrule
1 & DET & 159.0 & East \#5  & 1 & Neemias Queta (BOS)      & 282 & 76.0 \\
2 & HOU & 124.7 & West \#1  & 2 & Moussa Diabat\'e (CHA)   & 245 & 63.6 \\
3 & CLE & 123.1 & East \#2  & 3 & Ryan Kalkbrenner (CHA)   & 237 & 61.1 \\
4 & LAC & 121.9 & West \#12 & 4 & Jalen Duren (DET)        & 227 & 92.9 \\
5 & SAS & 121.7 & West \#3  & 5 & Kel'el Ware (MIL)        & 219 & 71.0 \\
6 & GSW & 117.0 & West \#7  & 6 & Ryan Rollins (MIL)       & 207 & 63.6 \\
7 & CHA & 114.6 & East \#11 & 7 & Payton Pritchard (BOS)   & 200 & 78.2 \\
8 & LAL & 112.3 & West \#6  & 8 & Cam Spencer (MEM)        & 197 & 50.7 \\
\bottomrule
\end{tabular}
\end{table}

Detroit is the standout: a projected playoff seed on a \$140M, tax-free
payroll, anchored by two projected All-NBA players of whom one (Duren) earns
\$9M. The six-season backtest gives that profile a precedent with both
outcomes: OKC 2023--24 converted the same shape into a title within two
seasons; Memphis 2021--23 topped the value board twice and never escaped the
second round.

\section{Limitations}
\begin{itemize}[itemsep=2pt]
\item \textbf{COVID seasons excluded.} 2019--20 and 2020--21 are omitted
because uneven schedules and pro-rated tax bills break the availability and
cost models respectively; the calibration window is therefore six seasons
spanning eight years, not eight consecutive seasons.
\item \textbf{Pre-2023 salary reconstruction.} Older seasons are rebuilt from
contract archives that purge superseded deals; coverage is restored via
individual salary histories and hand patches, and team totals are validated
against independent aggregates, but a handful of departed role players carry
imputed minimum salaries.
\item \textbf{Roster incompleteness for 2026--27.} Projections use contracts
signed as of the run date; unsigned free agents (e.g., LeBron James, James
Harden) and future signings are absent, which depresses their destination
teams. 2026--27 cap, tax, and apron figures are league projections.
\item \textbf{Six champions is still a small fit.} The title premiums
$(a,b) = (0.2, 1.8)$ carry wide uncertainty. Cross-validation
(Section~\ref{sec:cv}) shows the \emph{direction} is stable --- rotation over
star in all six folds --- but $b$ ranges from 1.8 to 2.4 across them, and
mean champion rank is the weakest of the three validated quantities.
\item \textbf{Descriptive tables use the full-sample fit.} The per-season
tables in Section~\ref{sec:backtest} and the 2026--27 projections are
generated with weights fit on all six seasons, which is the right choice for
description and forecasting. Only Section~\ref{sec:cv} reports accuracy the
model had not already seen; cite those figures, not the in-sample ones.
\end{itemize}

\vspace{6pt}
\hrule
\vspace{4pt}
{\small Generated from the R-Value pipeline
(\texttt{backtest.py}, \texttt{predict.py}); data: NBA Stats API and public
salary records, seasons 2018--19 through 2025--26 (ex-COVID), contracts as of
July 5, 2026. League average $=$ 100 for all R-Values.}

\end{document}
